% answers_short.tex —— 2011 年·湖北卷·理科 · 简短答案速查
% 由 tools/make_short_answers.py 自 answers.tex 的 \gksolhead/\gkanswer 提取，
% 并与 2011/2011_数学真题_湖北理_答案要点.md 逐题比对后生成，勿手改（改源 MD 后重跑）。

\begin{gknotice}
\textbf{使用说明}：本节只给最终答案便于快速核对，完整推导见解析版。答案为编写者独立推导，\textbf{非官方标准答案}。
\end{gknotice}

\gkpart{选择题}
\begin{center}\begin{tabular}{ccccccccccc}
\toprule
题号 & 1 & 2 & 3 & 4 & 5 & 6 & 7 & 8 & 9 & 10 \\
\midrule
答案 & A & A & B & C & C & B & B & D & C & D \\
\bottomrule
\end{tabular}\end{center}

\gkpart{填空题}
\begin{itemize}[leftmargin=2.6em]
  \item[\textbf{11.}] $17$。
  \item[\textbf{12.}] $\dfrac{28}{145}$。
  \item[\textbf{13.}] $\dfrac{67}{66}$。
  \item[\textbf{14.}] (1) $(2,2)$；(2) $(x-1)^2+y^2=1$
  \item[\textbf{15.}] $21$；$43$。
\end{itemize}

\gkpart{解答题}
\begin{itemize}[leftmargin=2.6em]
  \item[\textbf{16.}] (1) $5$；(2) $\dfrac{11}{16}$
  \item[\textbf{17.}] (1) $v(x)=\begin{cases}60,&0\leqslant x\leqslant20,\\[4pt]\dfrac{200-x}3,&20<x\leqslant200.\end{cases}$；(2) $x=100$ 时最大，约 $3333$ 辆/小时
  \item[\textbf{18.}] (1) 证明见解析；(2) $\tan\theta$ 的最小值为 $\dfrac{\sqrt6}3$
  \item[\textbf{19.}] (1) $a_1=a$，$a_n=ar(r+1)^{n-2}$（$n\geqslant2$）；(2) 是；由条件得 $r=-3$，此时 $a_{m+1},a_m,a_{m+2}$ 对任意 $m\geqslant2$ 成等差数列
  \item[\textbf{20.}] (1) $C:x^2-\dfrac{y^2}{m}=a^2$（$m\neq0$）；$m>0$ 为双曲线，$-1<m<0$ 为椭圆，$m=-1$ 为圆，$m<-1$ 为椭圆（焦点在 $y$ 轴）；(2) 存在条件为 $m\in\left[\dfrac{1-\sqrt5}2,0\right)\cup\left(0,\dfrac{1+\sqrt5}2\right]$；$m<0$ 时 $\tan\angle F_1NF_2=2$，$m>0$ 时 $\tan\angle F_1NF_2=-2$
  \item[\textbf{21.}] (1) $0$；(2) 证明见解析
\end{itemize}
